\chapter{Implémentation et validation}
\label{chap:implementation}

\section{Environnement technique}

\begin{itemize}
\item Langages : Python 3.10 (backend), TypeScript/React (frontend).
\item Frameworks : PyTorch, FastAPI, Streamlit (prototypage), React.
\item Librairies : MNE, Scikit-learn, SHAP, Captum.
\item Conteneurisation : Docker, Kubernetes (Minikube).
\end{itemize}

\section{Prétraitement des données}

Le jeu de données \textbf{Sleep-EDF} (version 2018) a été utilisé. Il contient 153 enregistrements PSG de sujets sains et insomniagues. Les étapes :
\begin{enumerate}
\item Chargement des fichiers EDF avec MNE.
\item Sélection des canaux EEG (Fpz-Cz, Pz-Oz) + optionnels EOG, EMG.
\item Filtrage passe-bande (0.5–35 Hz) et rééchantillonnage à 100 Hz.
\item Segmentation en épisodes de 30 secondes (3000 points).
\item Normalisation par épisode (z-score).
\item Division entraînement / validation / test (70/15/15) au niveau des sujets.
\end{enumerate}

\section{Entraînement du modèle d’ensemble}

Chaque modèle de base est entraîné séparément avec une validation croisée. Les hyperparamètres optimisés sont :

\begin{table}[H]
\centering
\caption{Hyperparamètres finaux}
\begin{tabular}{|l|l|}
\hline
Modèle & Paramètres \\
\hline
CNN & batch size 64, lr=1e-3, dropout 0.3 \\
LSTM & hidden size 128, 2 couches, lr=1e-3 \\
Transformer & 4 têtes, 128 dim, lr=1e-4 \\
Méta-apprenant & 2 couches (128,64), lr=1e-4, fenêtre LCS=5 \\
\hline
\end{tabular}
\end{table}

Le méta-apprenant est entraîné sur les prédictions des modèles de base (probabilités softmax). Le lissage contextuel est appliqué en post-traitement.

\section{Résultats expérimentaux}

\subsection{Performance de classification}

\begin{table}[H]
\centering
\caption{Précision par modèle (test)}
\begin{tabular}{|l|c|}
\hline
Modèle & Accuracy (\%) \\
\hline
CNN seul & 87.2 \\
LSTM seul & 92.4 \\
Transformer seul & 91.8 \\
Ensemble (stacking) sans LCS & 93.5 \\
\textbf{Ensemble avec LCS} & \textbf{94.1} \\
\hline
\end{tabular}
\end{table}

Le modèle final atteint une précision globale de 94,1\%, un score de Cohen's Kappa de 0,91, dépassant l’objectif des 93\%.

\subsection{Explicabilité : exemples}

\begin{figure}[H]
\centering
\fbox{\parbox{0.6\textwidth}{\centering \textbf{Exemple de valeurs SHAP pour un épisode REM}}}
\caption{Attribution SHAP pour un épisode correctement classé REM}
\label{fig:shap}
\end{figure}

La figure~\ref{fig:shap} montre que les échantillons autour de 250-500 ms (mouvements oculaires) contribuent positivement à la classe REM. Les cartes de saillie du CNN confirment l’importance de la région occipitale (canal Pz-Oz) pour la détection des ondes lentes.

\section{Validation par des experts}

Un petit panel de trois techniciens du sommeil a évalué 100 épisodes aléatoires avec les explications fournies. Le taux d’accord avec les prédictions a augmenté de 85\% (sans explications) à 92\% (avec explications), démontrant l’utilité clinique de la transparence.

\section{Déploiement}

L’application a été conteneurisée et déployée sur un cluster Kubernetes local. Les temps de réponse moyens sont de 1,2 seconde par épisode (inférence complète sur une nuit de 8h en 8 minutes), respectant les exigences.